% Efficient attention study zone.

\section*{FlashAttention 与 FlexAttention 学习专区}
\phantomsection
\addcontentsline{toc}{section}{FlashAttention 与 FlexAttention 学习专区}

\begin{conceptbox}
\textbf{学习目标：} 从 AIGC 工程使用者视角掌握高效注意力：\key{知道原理直觉、会启用和选择实现、能读懂别人代码里的 attention backend / mask / kernel 相关配置}，而不是深入研究 CUDA kernel 或 FlashAttention 论文细节。\\
\textbf{当前主线：} 按 \key{标准 attention 瓶颈 $\rightarrow$ FlashAttention 核心直觉 $\rightarrow$ PyTorch SDPA / flash-attn 使用入口 $\rightarrow$ 常见代码模式 $\rightarrow$ FlexAttention 自定义 mask $\rightarrow$ AIGC 场景选型与调试} 的顺序推进。
\end{conceptbox}

% ============================================================
% FlashAttention 与 FlexAttention 学习路线图
% ============================================================

\section{学习路线图：FlashAttention 与 FlexAttention}

\begin{conceptbox}
\textbf{一句话目标：} 你不是要研究 FlashAttention 或 FlexAttention 本身，而是要成为一个\key{会用、看得懂、能判断 trade-off 的 AIGC 工程使用者}。学完后你应该能回答四件事：\\
1. \key{标准 attention 为什么在长序列、多图像 token、视频 token 场景下又慢又吃显存}；\\
2. \key{FlashAttention 大概怎么做到 exact、快、省显存}，但不需要深入 CUDA kernel 细节；\\
3. \key{看到 PyTorch、Transformers、Diffusers、vLLM 里的 attention backend 配置时，知道代码在做什么}；\\
4. \key{什么时候用普通 SDPA / FlashAttention，什么时候才需要 FlexAttention 这种可编程 mask 抽象。}
\end{conceptbox}

\subsection*{资料收集与引用口径}

\begin{conceptbox}
\textbf{本专区先查资料再整理结论。} 这一页主要综合五类资料：\\
1. \key{FlashAttention 系列论文}：FlashAttention、FlashAttention-2、FlashAttention-3，用来理解 IO-aware exact attention、并行划分和新硬件优化的来源；\\
2. \key{官方实现资料}：Dao-AILab \texttt{flash-attention} 仓库，用来确认工程接口、支持硬件、安装和常见使用方式；\\
3. \key{PyTorch 官方资料}：\texttt{scaled\_dot\_product\_attention} 与 \texttt{torch.nn.attention.flex\_attention}，用来理解 PyTorch 内置 attention backend 与 FlexAttention API；\\
4. \key{PyTorch 官方博客}：FlexAttention blog，用来理解它为什么强调“PyTorch 的灵活性 + FlashAttention 的性能”；\\
5. \key{已有系统笔记}：显存账本、mixed precision、activation checkpointing 与 KV cache，用来把 attention kernel 优化放回训练/推理系统上下文里。
\end{conceptbox}

\subsection{一、这个专区的学习定位}

\begin{warnbox}
\textbf{不要把目标学偏：} 你当前的目标不是复现 FlashAttention kernel，也不是跟进每一个 CUDA/Triton 优化细节；你的目标是：\key{知道它为什么有用、实际项目里怎么启用、读别人代码时能识别它、出问题时知道排查方向。}
\end{warnbox}

从 AIGC 工程角度看，FlashAttention / FlexAttention 通常会出现在这些地方：
\begin{itemize}[leftmargin=1.5em]
  \item 训练 LLM、DiT、视频模型时，想降低 attention 激活显存和提升吞吐。
  \item 使用 Hugging Face \texttt{Transformers} 时看到 \texttt{attn\_implementation=\"flash\_attention\_2\"}、\texttt{sdpa}、\texttt{eager}。
  \item 使用 PyTorch 2.x 时看到 \texttt{scaled\_dot\_product\_attention}，但不知道它底层选了哪个 backend。
  \item 看推理框架或训练框架时看到 \texttt{flash\_attn\_varlen\_func}、\texttt{cu\_seqlens}、\texttt{causal}、\texttt{sliding\_window}。
  \item 做长上下文、多图、多帧视频或文档边界 mask 时，发现普通 dense mask 太慢或太占显存，开始接触 FlexAttention。
\end{itemize}

\begin{summarybox}
\textbf{本专区的主线：}\\
不是“论文作者怎么设计 kernel”，而是 \key{普通 attention 为什么不够用 $\rightarrow$ FlashAttention 为什么能作为默认优化 $\rightarrow$ 代码里怎么打开和识别 $\rightarrow$ FlexAttention 什么时候有必要 $\rightarrow$ AIGC 项目怎么选型和排查。}
\end{summarybox}

\subsection{二、先建立位置关系：它们各自解决什么}

\begin{center}
{\small
\begin{tabular}{p{3.2cm}p{4.4cm}p{5.3cm}}
\hline
\textbf{主题} & \textbf{你需要掌握的程度} & \textbf{使用者视角的一句话理解} \\
\hline
标准 Attention & 必须熟悉公式与瓶颈 & 每个 token 看所有 token，表达力强，但长序列下 $n^2$ 中间矩阵和 IO 很贵。 \\
PyTorch SDPA & 必须会识别和使用 & PyTorch 2.x 的统一 attention 入口，底层可能走 math、memory-efficient 或 FlashAttention 类 backend。 \\
FlashAttention & 必须懂核心直觉与使用方式 & 不改 attention 数学结果，通过分块和 online softmax 避免完整物化 attention matrix，通常是标准 dense/causal attention 的优先优化。 \\
FlashAttention-2/3 & 了解版本差异即可 & 主要是更好的并行划分和新硬件利用率；作为使用者通常只需要知道版本、硬件、dtype 和框架兼容性。 \\
\texttt{flash-attn} 包 & 会看常见接口即可 & 在很多训练/推理项目里直接调用，常见关键词是 varlen、cu\_seqlens、causal、dropout、head dim。 \\
FlexAttention & 会判断适用场景 & 不是替代所有 FlashAttention，而是用于复杂 score/mask/block-sparse 需求下的可编程高性能 attention。 \\
\hline
\end{tabular}
}
\end{center}

\begin{intubox}
\textbf{最重要的区分：} FlashAttention 解决的是\key{标准 attention 的实现效率}；FlexAttention 解决的是\key{复杂 attention pattern 的表达与性能}。\\
如果只是普通 causal LLM 或 dense DiT attention，优先考虑 SDPA / FlashAttention；如果有 prefix-LM、sliding window、document mask、block-sparse、图文分区 mask，再考虑 FlexAttention。
\end{intubox}

\subsection{三、正式学习顺序（工程使用者版）}

\subsubsection*{Stage 0：先知道普通 Attention 为什么贵}

\begin{itemize}[leftmargin=1.5em]
  \item \textbf{学什么：} scaled dot-product attention 公式、$QK^\top$、softmax、mask、dropout、$PV$，以及 $n \times n$ score/probability matrix 为什么导致显存和 IO 压力。
  \item \textbf{使用者为什么要懂：} 你以后判断 OOM 或训练慢时，不能只说“开 FlashAttention”，而要知道它主要优化的是 attention 的中间矩阵物化和内存访问。
  \item \textbf{完成标志：} 看到一个长序列或视频 token 数量，你能直觉判断 attention 可能成为瓶颈。
\end{itemize}

\subsubsection*{Stage 1：掌握 FlashAttention 的核心直觉，不钻 kernel}

\begin{itemize}[leftmargin=1.5em]
  \item \textbf{学什么：} exact attention、tiling、online softmax、少存中间矩阵、backward 重算一部分中间量、HBM/SRAM IO 差异。
  \item \textbf{不需要学到什么：} 不需要逐行读 CUDA/Triton kernel，不需要掌握 warp-level scheduling，不需要复现论文里的 IO complexity 证明。
  \item \textbf{完成标志：} 你能说清楚：FlashAttention \key{不是近似 attention，不是 sparse attention，也不是改模型结构}；它主要是更聪明地算同一个 attention。
\end{itemize}

\subsubsection*{Stage 2：学会项目里怎么启用与识别 FlashAttention / SDPA}

\begin{itemize}[leftmargin=1.5em]
  \item \textbf{学什么：} \texttt{torch.nn.functional.scaled\_dot\_product\_attention}、PyTorch SDPA backend、Hugging Face \texttt{attn\_implementation}、\texttt{flash\_attention\_2}、\texttt{sdpa}、\texttt{eager}，以及 \texttt{flash-attn} 包的常见接口。
  \item \textbf{为什么这是重点：} 你大多数时候不是手写 attention kernel，而是在别人的模型代码、训练脚本、配置文件里看到这些开关。
  \item \textbf{完成标志：} 看到代码里使用 \texttt{attn\_implementation=\"flash\_attention\_2\"} 或 \texttt{scaled\_dot\_product\_attention}，你知道它大概影响性能、显存、兼容性和 mask 支持。
\end{itemize}

\subsubsection*{Stage 3：读懂常见代码模式和坑}

\begin{itemize}[leftmargin=1.5em]
  \item \textbf{学什么：} \texttt{causal} mask、padding mask、varlen attention、\texttt{cu\_seqlens}、\texttt{max\_seqlen}、head dimension 限制、dtype 限制、dropout、training/eval 行为差异。
  \item \textbf{为什么要学：} 很多 bug 不是“不懂论文”，而是 mask 形状错、padding 没处理、版本不兼容、head dim 不支持、dtype 不满足 kernel 条件。
  \item \textbf{完成标志：} 看别人项目里的 attention forward，你能判断 Q/K/V shape、mask 类型、是否 varlen、是否 causal，以及为什么要这样写。
\end{itemize}

\subsubsection*{Stage 4：再学 FlexAttention，重点放在“什么时候需要”}

\begin{itemize}[leftmargin=1.5em]
  \item \textbf{学什么：} \texttt{score\_mod}、\texttt{mask\_mod}、\texttt{BlockMask}、\texttt{flex\_attention}、\texttt{torch.compile}，以及 sliding window、prefix-LM、document mask、block-sparse attention 的表达方式。
  \item \textbf{学习重点：} FlexAttention 的价值不是“比 FlashAttention 更高级”，而是当 attention pattern 不是普通 dense/causal 时，还想保留接近 fused kernel 的性能。
  \item \textbf{完成标志：} 给你一个业务需求，例如“不同文档之间不能互相 attend”或“视频只看局部窗口 + 少量全局 token”，你能判断是否适合 FlexAttention。
\end{itemize}

\subsubsection*{Stage 5：放回 AIGC 场景做选型}

\begin{itemize}[leftmargin=1.5em]
  \item \textbf{学什么：} LLM 长上下文、DiT 高分辨率图像 token、视频生成、多图多文本条件、KV cache、GQA/MQA、activation checkpointing 与 FlashAttention 的组合。
  \item \textbf{为什么最后学：} 真实项目里你关心的是“我该开哪个实现、为什么 OOM、为什么速度没变快、mask 支不支持”，不是单独背一个 kernel 名。
  \item \textbf{完成标志：} 训练 OOM、推理慢、复杂 mask、长上下文四类问题出现时，你能给出 attention 层面的排查顺序。
\end{itemize}

\subsection{四、建议固化的笔记树}

\begin{itemize}[leftmargin=1.5em]
  \item \texttt{00\_learning\_roadmap.tex}：这页使用者路线图。
  \item \texttt{01\_attention\_bottleneck\_ledger.tex}：标准 attention 的计算、显存与 IO 账本。
  \item \texttt{02\_flashattention\_core\_intuition.tex}：FlashAttention 的 exact attention、tiling、online softmax 直觉。
  \item \texttt{03\_sdpa\_and\_flashattn\_usage.tex}：PyTorch SDPA、Hugging Face \texttt{attn\_implementation}、\texttt{flash-attn} 包常见入口。
  \item \texttt{04\_attention\_code\_reading\_patterns.tex}：Q/K/V shape、causal mask、padding mask、varlen、\texttt{cu\_seqlens}、dtype/head dim 限制。
  \item \texttt{05\_flexattention\_for\_custom\_masks.tex}：FlexAttention 的 \texttt{score\_mod}、\texttt{mask\_mod}、\texttt{BlockMask} 与适用场景。
  \item \texttt{06\_aigc\_attention\_selection\_and\_debugging.tex}：LLM、DiT、视频生成、长上下文和复杂 mask 的选型与排查总结。
\end{itemize}

\subsection{五、你现在最应该掌握到什么深度？}

\begin{summarybox}
\textbf{必须掌握：} 标准 attention 的 $O(n^2)$ 瓶颈；FlashAttention 为什么 exact 但省显存；PyTorch SDPA / Hugging Face / \texttt{flash-attn} 常见使用入口；常见 mask、shape、dtype、head dim 限制；FlexAttention 适合什么复杂 mask。\\
\textbf{概念了解：} FlashAttention-2/3 相比 FlashAttention 1 大致优化了并行度、硬件利用率和低精度路径；知道版本和硬件兼容性会影响实际效果。\\
\textbf{可以暂时不学：} CUDA kernel 源码、Triton kernel 编写、warp/block 级调度、Hopper WGMMA/TMA 指令、论文中完整 IO complexity 证明。
\end{summarybox}

\subsection{六、如果你现在就开始，我建议的第一轮学习节奏}

\begin{summarybox}
\textbf{最优起步顺序：}\\
1. 先学标准 attention 的显存账本，知道 $n^2$ 中间矩阵和 IO 压力从哪来；\\
2. 再学 FlashAttention 核心直觉，重点是 exact、tiling、online softmax、少存中间矩阵；\\
3. 立刻进入使用入口：PyTorch SDPA、Hugging Face \texttt{attn\_implementation}、\texttt{flash-attn} 常见接口；\\
4. 然后专门学一页代码阅读模式，解决“看到别人代码能看懂”的问题；\\
5. 最后学 FlexAttention，重点是复杂 mask 什么时候值得用。\\
换句话说，\key{先懂瓶颈和直觉，再学接口和代码，最后再学可编程复杂 mask。}
\end{summarybox}

\subsection{七、阶段性面试 / 代码评审版总结}

如果面试官或同事问“FlashAttention 和 FlexAttention 你怎么理解”，可以这样答：

\begin{conceptbox}
FlashAttention 是一种 \key{高性能 exact attention 实现}。它不改变 attention 的数学结果，而是通过 tiling 和 online softmax 避免完整物化 attention matrix，减少 HBM 读写和训练显存峰值。对使用者来说，它通常是标准 dense/causal attention 的默认优化，重点是看框架是否启用、dtype/head dim/mask 是否支持、实际是否命中高性能 backend。FlexAttention 则更偏 \key{可编程 attention 抽象}，适合 sliding window、prefix-LM、document mask、block-sparse 等复杂模式；它不是替代所有 FlashAttention，而是在复杂 mask 场景下兼顾灵活性和性能。
\end{conceptbox}

% ============================================================
\subsection{参考资料}
% ============================================================

\begingroup
\small
\sloppy
\begin{itemize}[leftmargin=1.5em]
  \item Dao 等，\emph{FlashAttention: Fast and Memory-Efficient Exact Attention with IO-Awareness}, arXiv:2205.14135。\\
  \url{https://arxiv.org/abs/2205.14135}
  \item Dao，\emph{FlashAttention-2: Faster Attention with Better Parallelism and Work Partitioning}, arXiv:2307.08691。\\
  \url{https://arxiv.org/abs/2307.08691}
  \item Shah 等，\emph{FlashAttention-3: Fast and Accurate Attention with Asynchrony and Low-precision}, arXiv:2407.08608。\\
  \url{https://arxiv.org/abs/2407.08608}
  \item Dao-AILab \texttt{flash-attention} 官方仓库：FlashAttention 安装、接口、支持硬件与使用说明。\\
  \url{https://github.com/Dao-AILab/flash-attention}
  \item PyTorch \texttt{scaled\_dot\_product\_attention} 文档：PyTorch 内置 SDPA API 与 backend 语义。\\
  \url{https://docs.pytorch.org/docs/stable/generated/torch.nn.functional.scaled_dot_product_attention.html}
  \item PyTorch FlexAttention 官方博客：\emph{FlexAttention: The Flexibility of PyTorch with the Performance of FlashAttention}。\\
  \url{https://pytorch.org/blog/flexattention/}
  \item PyTorch \texttt{torch.nn.attention.flex\_attention} 文档：\texttt{flex\_attention}、\texttt{score\_mod}、\texttt{BlockMask} 等 API。\\
  \url{https://docs.pytorch.org/docs/stable/nn.attention.flex_attention.html}
\end{itemize}
\endgroup

% ============================================================
% Stage 0：标准 Attention 的瓶颈账本
% ============================================================

\section{Stage 0：标准 Attention 的瓶颈账本}

\begin{conceptbox}
\textbf{这一节的目标：} 先不急着学 FlashAttention。你要先把标准 attention 到底贵在哪里讲清楚。学完这一节，你至少要能回答四件事：\\
1. \key{标准 scaled dot-product attention 的每一步在算什么}；\\
2. \key{$O(n^2)$ 到底来自哪里，是计算量、显存，还是 IO}；\\
3. \key{训练和推理阶段 attention 的显存瓶颈有什么不同}；\\
4. \key{为什么 FlashAttention 能成为 AIGC 工程里的默认优化方向。}
\end{conceptbox}

\note{这一节只建立瓶颈账本，不深入 FlashAttention kernel。你可以把它理解成后面所有高效 attention 技术的“问题定义”。}

\subsection{一、先从标准 attention 公式开始}

Transformer 里的 scaled dot-product attention 通常写成：

\[
\mathrm{Attention}(Q,K,V)
=
\mathrm{softmax}\left(\frac{QK^\top}{\sqrt{d}} + M\right)V
\]

其中：
\begin{itemize}[leftmargin=1.5em]
  \item $Q$ 是 query，表示“当前位置想找什么信息”。
  \item $K$ 是 key，表示“每个位置能被什么方式匹配”。
  \item $V$ 是 value，表示“匹配之后真正要取回的信息”。
  \item $d$ 是每个 head 的维度，$\sqrt{d}$ 用来缩放 dot product，避免 logits 过大导致 softmax 饱和。
  \item $M$ 是 mask，例如 causal mask、padding mask、局部窗口 mask、文档边界 mask。
\end{itemize}

从工程实现上看，单个 attention head 的核心流程可以拆成四步：

\begin{center}
{\small
\begin{tabular}{p{3.2cm}p{4.2cm}p{5.4cm}}
\hline
\textbf{步骤} & \textbf{张量形状} & \textbf{含义} \\
\hline
输入 $Q,K,V$ & $Q,K,V \in \mathbb{R}^{n \times d}$ & $n$ 是 token 数，$d$ 是 head dim。 \\
算 score & $S = QK^\top \in \mathbb{R}^{n \times n}$ & 每个 query token 和每个 key token 两两打分。 \\
做 softmax & $P = \mathrm{softmax}(S / \sqrt{d} + M)$ & 把分数变成每个 query 对所有 key 的注意力分布。 \\
加权求和 & $O = PV \in \mathbb{R}^{n \times d}$ & 用注意力分布从 value 里取信息。 \\
\hline
\end{tabular}
}
\end{center}

\begin{intubox}
\textbf{最关键的直觉：} 标准 attention 的语义是\key{每个 token 都和所有 token 发生一次关系判断}。所以只要 token 数 $n$ 变大，token pair 数量就会按 $n^2$ 增长。
\end{intubox}

\subsection{二、$O(n^2)$ 到底在哪里出现？}

最容易混淆的一点是：$O(n^2)$ 不只是“算得多”，它还意味着中间矩阵可能很大，而且这些中间矩阵还要在 GPU 显存层级之间反复读写。

\subsubsection*{1. 计算量：每个 query 都要扫所有 key}

单个 head 的 score 计算是：

\[
S = QK^\top
\]

$Q$ 的形状是 $n \times d$，$K^\top$ 的形状是 $d \times n$，所以输出 $S$ 是 $n \times n$。这个矩阵里每一个元素 $S_{ij}$ 都表示：第 $i$ 个 query token 和第 $j$ 个 key token 的相似度。

因此 score matmul 的计算量大约是：

\[
O(n^2 d)
\]

后面的 $PV$ 也是类似量级：

\[
P \in \mathbb{R}^{n \times n},\quad V \in \mathbb{R}^{n \times d},\quad O=PV \Rightarrow O(n^2 d)
\]

所以 attention 的主要矩阵乘计算可以粗略理解成两次 $n^2d$ 级别的操作。

\subsubsection*{2. 显存：score / probability matrix 是 $n \times n$}

如果朴素实现把 score matrix $S$ 和 softmax 后的 probability matrix $P$ 都显式保存下来，那么单个 batch、单个 head 就会有 $n^2$ 个元素。

如果考虑 batch size $B$ 和 attention heads $H$，那么 attention probability 的元素数量就是：

\[
B \times H \times n \times n
\]

若 dtype 是 BF16/FP16，每个元素 2 bytes，那么只保存一个 attention matrix 的字节数约为：

\[
M_{\text{attn matrix}}
\approx
B \times H \times n^2 \times 2
\]

\begin{warnbox}
\textbf{注意：} 这还只是一个矩阵的大小。训练时为了反向传播，朴素实现可能还需要保存 softmax 结果、dropout mask、部分中间状态，真实峰值会更复杂。
\end{warnbox}

\subsubsection*{3. 先把 HBM / SRAM 讲清楚：IO 不是磁盘 IO，而是 GPU 内存搬运}

这里的 \key{IO} 不是读文件、读网络，也不是 CPU 磁盘 IO。FlashAttention 论文里说的 IO，主要是指 \key{GPU 上不同层级内存之间的数据搬运}，尤其是 HBM 和片上高速内存之间的读写。

\begin{conceptbox}
\textbf{一句话版：} HBM 是 GPU 的“大仓库”，容量大，能长期放 tensor，但访问它仍然比片上内存贵；SRAM 是 GPU 芯片里的“小工作台”，非常快，但容量很小，只能临时放 tile、局部变量和中间结果。FlashAttention 的核心直觉就是：\key{少把大矩阵搬进搬出大仓库，多在小工作台上分块算完就合并结果。}
\end{conceptbox}

先按工程直觉理解这几个层级：

\begin{center}
{\small
\begin{tabular}{p{2.9cm}p{3.2cm}p{3.4cm}p{3.5cm}}
\hline
\textbf{层级} & \textbf{大概是什么} & \textbf{特点} & \textbf{在 attention 里的角色} \\
\hline
HBM / global memory & GPU 显存主体，CUDA tensor 大多长期存这里 & 容量大，带宽高，但相对片上内存仍然慢 & 放 $Q,K,V,O$，朴素实现还会放完整 $S$ 和 $P$。 \\
L2 cache & GPU 芯片上的共享缓存 & 比 HBM 快，容量远小于 HBM，由硬件自动管理 & 缓存最近访问的数据，提升复用，但不能指望它装下完整 attention matrix。 \\
Shared memory / L1 & 每个 SM 附近的片上高速存储 & 很快，容量很小，常由 kernel 显式组织使用 & 放当前 block 正在处理的 $Q/K/V$ tile。 \\
Registers & 每个线程自己的寄存器 & 最快，但最小 & 放 running max、running sum、局部累加结果等标量或小向量。 \\
\hline
\end{tabular}
}
\end{center}

\begin{intubox}
\textbf{怎么理解论文里的 SRAM：} FlashAttention 论文把 HBM 和 SRAM 抽象成两级内存模型。这里的 SRAM 不是说你在代码里一定会看到一个变量叫 \texttt{SRAM}，而是泛指 GPU 芯片上的高速小容量存储，例如 registers、shared memory、L1/L2 cache 等。实际 CUDA/Triton kernel 会综合使用这些层级。
\end{intubox}

\subsubsection*{4. 为什么不能直接把 attention matrix 放进 SRAM？}

因为 attention matrix 太大，而 SRAM 太小。以前面这个例子看：

\[
B=1,\quad H=32,\quad n=8192,\quad \text{dtype}=\mathrm{BF16/FP16}
\]

完整 attention matrix 的大小是：

\[
1 \times 32 \times 8192^2 \times 2
\approx
4.0\ \mathrm{GiB}
\]

如果只看\key{单个 head}，也有：

\[
8192^2 \times 2
= 128\ \mathrm{MiB}
\]

而一个 GPU block 能直接高效使用的 shared memory / registers 只是很小的一块片上工作区，远远装不下一个完整的 $8192 \times 8192$ 矩阵。

\begin{warnbox}
\textbf{所以 FlashAttention 不是“把完整 $S$ 放到 SRAM 里”。}\\
真正做法是：每次只拿一小块 $Q$、一小块 $K/V$ 到片上高速内存里，算出当前 tile 的局部 score；用 online softmax 把局部结果合并到最终输出；然后丢掉这个局部 score tile，继续处理下一块。这样就避免把完整 $S$ 或 $P$ 写回 HBM。
\end{warnbox}

\subsubsection*{5. 朴素 attention 的 HBM 往返路径}

朴素 attention 的问题不是只算了一次 $S$，而是可能经历这样的路径：

\begin{center}
\begin{tabular}{p{13.2cm}}
\hline
从 HBM 读 $Q,K$，计算 $S=QK^\top$，把完整 $S$ 写回 HBM \\
$\downarrow$ \\
再从 HBM 读 $S$，做 mask 和 softmax，把完整 $P$ 写回 HBM \\
$\downarrow$ \\
再从 HBM 读 $P,V$，计算 $O=PV$，把 $O$ 写回 HBM \\
\hline
\end{tabular}
\end{center}

如果 $S$ 和 $P$ 都是 $B \times H \times n \times n$ 级别，那么这里最浪费的不是“有一次矩阵乘”，而是\key{把巨大的中间矩阵写到 HBM，再从 HBM 读回来}。

\subsubsection*{6. FlashAttention 的 IO-aware 到底是什么意思？}

\key{IO-aware} 的意思是：设计 attention 算法时，不只数 FLOPs，也要数 HBM 和片上内存之间搬了多少数据。

它的思路可以粗略理解成：

\begin{center}
\begin{tabular}{p{13.2cm}}
\hline
从 HBM 读一块 $Q$，再分批读 $K,V$ 的 tile 到片上高速内存 \\
$\downarrow$ \\
在 tile 内部算局部 score，并用 online softmax 更新 running max / running sum \\
$\downarrow$ \\
不保存完整 $S/P$，只维护当前输出块的累加结果 \\
$\downarrow$ \\
最后把输出 $O$ 写回 HBM \\
\hline
\end{tabular}
\end{center}

\begin{intubox}
\textbf{FlashAttention 论文最重要的切入点就是这里：} 标准 attention 的数学公式没有错，瓶颈在于朴素实现会把巨大的 $n \times n$ 中间矩阵写回 HBM，再读出来继续算。FlashAttention 要优化的正是这种\key{HBM 读写次数、内存访问路径和中间矩阵物化}。
\end{intubox}

\begin{summarybox}
\textbf{你现在应该形成的准确理解：} HBM 约等于 GPU 上放大 tensor 的主显存；SRAM 约等于 GPU 芯片上很快但很小的工作区。FlashAttention 不是把全部 attention matrix 塞进 SRAM，而是把计算切成很多 tile，在 SRAM 级别的小工作区里一块块处理，并用 online softmax 合并结果，从而避免完整物化 $S/P$。
\end{summarybox}

\subsubsection*{7. online softmax：为什么分块以后还能等价于标准 softmax？}

上面一直说 FlashAttention 可以一块一块处理 score tile，并用 online softmax 合并结果。这里先把 online softmax 单独讲清楚。

\begin{conceptbox}
\textbf{一句话版：} online softmax 是一种\key{边读 logits、边维护 softmax 所需统计量}的方法。它不要求一次性拿到完整 logits 向量，也不需要先把整行 $S$ 存下来；只要维护 running max、running sum 和输出累加器，就能得到和普通 stable softmax 等价的结果。
\end{conceptbox}

普通 softmax 对一行 logits $x=[x_1,x_2,\dots,x_n]$ 的定义是：

\[
\mathrm{softmax}(x_i)=\frac{e^{x_i}}{\sum_j e^{x_j}}
\]

工程里一般不会直接这样算，因为 $e^{x_i}$ 可能数值溢出。更稳定的写法是先减去全局最大值：

\[
m=\max_j x_j
\]

\[
\mathrm{softmax}(x_i)=\frac{e^{x_i-m}}{\sum_j e^{x_j-m}}
\]

这个写法的麻烦在于：朴素实现需要先看到整行 logits，才能知道全局最大值 $m$ 和分母。

online softmax 的思路是：如果 logits 是一块一块来的，也可以边处理边维护两个量：

\[
m=\text{目前看过的最大值}
\]

\[
l=\sum_{j \in \text{已看过的位置}} e^{x_j-m}
\]

这里的 $l$ 可以理解成当前 softmax 分母，只不过它永远是基于当前最大值 $m$ 归一化后的分母。

如果已经处理过一部分 logits，状态是 $m_{old},l_{old}$，现在新来一个 logit $x$，新的最大值是：

\[
m_{new}=\max(m_{old},x)
\]

旧分母 $l_{old}$ 是基于旧最大值 $m_{old}$ 算的；如果最大值变了，就必须把旧分母缩放到新最大值的坐标系下：

\[
l_{new}
=
e^{m_{old}-m_{new}}l_{old}
+
e^{x-m_{new}}
\]

\begin{intubox}
\textbf{关键直觉：} 如果后面来了一个更大的 logit，那么以前所有 $e^{x_j-m_{old}}$ 都要按照 $e^{m_{old}-m_{new}}$ 缩小。这样前后所有项都被统一到新的最大值 $m_{new}$ 下面，数值稳定性仍然成立。
\end{intubox}

FlashAttention 里通常不是一个 logit 一个 logit 地更新，而是一个 tile 一个 tile 地更新。假设新来一块 logits $x_{block}$，先算这一块自己的最大值和分母：

\[
m_{block}=\max(x_{block})
\]

\[
l_{block}=\sum_{x_i\in x_{block}}e^{x_i-m_{block}}
\]

再和旧状态合并：

\[
m_{new}=\max(m_{old},m_{block})
\]

\[
l_{new}
=
e^{m_{old}-m_{new}}l_{old}
+
e^{m_{block}-m_{new}}l_{block}
\]

这说明即使没有一次性保存完整 logits，也能得到和普通 stable softmax 相同的全局最大值和全局分母。

放回 attention 里，我们最终要的不是单独的 softmax probability matrix，而是：

\[
O=\mathrm{softmax}(S)V
\]

所以 FlashAttention 还会维护一个输出累加器：

\[
o=\sum_j e^{x_j-m}v_j
\]

当新 block 到来时，输出累加器也要和分母一样缩放到新的最大值坐标系：

\[
o_{new}
=
e^{m_{old}-m_{new}}o_{old}
+
e^{m_{block}-m_{new}}o_{block}
\]

最后再做归一化：

\[
O=\frac{o}{l}
\]

\begin{warnbox}
\textbf{online softmax 不是近似算法。} 它没有改变 softmax 的数学定义，也没有少看 token。它只是把“先存完整 logits 再统一 softmax”的流程，改成了“边分块扫描、边更新 running max / running sum / output accumulator”。在浮点误差范围内，它仍然是在算标准 softmax attention。
\end{warnbox}

\begin{summarybox}
\textbf{看代码时的理解方式：} 如果你看到 FlashAttention 或类似 kernel 里维护 \texttt{m}、\texttt{l}、\texttt{acc}、\texttt{rowmax}、\texttt{rowsum} 这类变量，它们通常就在做 online softmax：\texttt{m} 负责数值稳定的 running max，\texttt{l} 负责 softmax denominator，\texttt{acc} 负责累加 $\mathrm{softmax}(S)V$ 的未归一化输出。
\end{summarybox}

\subsection{三、一个长序列显存口算}

现在做一个最小口算，帮助你建立直觉。假设：

\begin{itemize}[leftmargin=1.5em]
  \item batch size $B=1$；
  \item attention heads $H=32$；
  \item sequence length $n=8192$；
  \item attention matrix 使用 BF16/FP16，2 bytes / element。
\end{itemize}

那么只保存一个 $B \times H \times n \times n$ attention matrix 就需要：

\[
1 \times 32 \times 8192^2 \times 2
\approx
4.29 \times 10^9 \text{ bytes}
\approx
4.0 \text{ GiB}
\]

这只是\key{一层 attention 的一个大中间矩阵}。如果训练时每层都需要保存相关中间状态，或者 batch 更大、序列更长、head 更多，显存压力会非常快地爆炸。

再看 $n=32768$：

\[
1 \times 32 \times 32768^2 \times 2
\approx
68.7 \times 10^9 \text{ bytes}
\approx
64.0 \text{ GiB}
\]

\begin{warnbox}
\textbf{这就是长上下文为什么难：} sequence length 从 8k 增到 32k，是变成 4 倍；但 attention matrix 是 $n^2$，所以相关中间矩阵会变成 16 倍。
\end{warnbox}

\subsection{四、AIGC 里为什么 attention 瓶颈特别常见？}

你是 AIGC 方向，所以不要只把 $n$ 理解成文本 token。很多生成模型都会把不同模态变成 token 序列：

\begin{center}
{\small
\begin{tabular}{p{3cm}p{4.4cm}p{5.5cm}}
\hline
\textbf{场景} & \textbf{$n$ 来自哪里} & \textbf{为什么容易变大} \\
\hline
LLM 长上下文 & 文本 token 数 & 从 4k、8k 到 32k、128k，$n^2$ 压力迅速上升。 \\
Diffusion Transformer / DiT & 图像 patch token & 分辨率越高，patch 数越多；多尺度或高分辨率训练更贵。 \\
视频生成 & 帧数 $\times$ 空间 patch token & 时间维度和空间维度一起放大，token 数很容易爆炸。 \\
多图多文本生成 & 文本 token + 图像 token & 条件信息更多，cross/self attention 的 token 组合更复杂。 \\
多文档 / RAG / 长 prompt & 多段文档 token & 还可能需要 document mask、prefix mask、局部窗口等复杂 mask。 \\
\hline
\end{tabular}
}
\end{center}

\begin{summarybox}
\textbf{AIGC 使用者视角：} Attention 优化不是一个“只在 LLM 里有用”的技巧。只要你的模型把文本、图像 patch、视频帧、条件信息都组织成 token，$n$ 变大以后 attention 就会变成核心瓶颈之一。
\end{summarybox}

\subsection{五、训练和推理的 attention 显存瓶颈不同}

\subsubsection*{1. 训练：主要怕激活和中间状态}

训练时要做 backward，所以 forward 阶段的很多中间状态要么需要保存，要么需要在 backward 时重算。对 attention 来说，朴素实现中最典型的压力来自：

\begin{itemize}[leftmargin=1.5em]
  \item $Q,K,V$ projection 输出；
  \item attention score matrix $S$；
  \item softmax probability matrix $P$；
  \item dropout mask；
  \item attention output 以及后续 residual / MLP 的激活。
\end{itemize}

因此训练 OOM 时，attention 相关优化通常会和这些策略一起出现：

\begin{itemize}[leftmargin=1.5em]
  \item \textbf{FlashAttention：} 减少 $S/P$ 这类 $n^2$ 中间矩阵的显式保存和 HBM 往返。
  \item \textbf{Activation checkpointing：} 少保存部分激活，反向时重算。
  \item \textbf{Mixed precision：} 用 BF16/FP16 降低参数、激活和中间状态的字节数。
  \item \textbf{Sequence parallel / tensor parallel：} 把某些张量在多个 GPU 上切分。
\end{itemize}

\subsubsection*{2. 推理：prefill 和 decode 是两种瓶颈}

LLM 推理通常可以分成两个阶段：

\begin{center}
{\small
\begin{tabular}{p{3cm}p{4.6cm}p{5.3cm}}
\hline
\textbf{阶段} & \textbf{在做什么} & \textbf{attention 瓶颈} \\
\hline
Prefill & 一次性处理整个 prompt & prompt 长时仍然接近训练里的 dense causal attention，容易受 $n^2$ 影响。 \\
Decode & 每次生成一个或少量新 token & 新 token 查询历史 KV cache，单步计算近似随历史长度线性增长。 \\
\hline
\end{tabular}
}
\end{center}

推理时通常不会像训练那样保存完整 backward 所需激活，但会长期保存 KV cache：

\[
M_{\text{KV cache}}
\propto
B \times L \times n \times H_{kv} \times d \times 2 \times \text{bytes}
\]

其中 $L$ 是层数，$H_{kv}$ 是 KV heads 数，$d$ 是 head dim，额外的 $2$ 表示 K 和 V 两份缓存。

\begin{intubox}
\textbf{一句话区分：} 训练更怕\key{attention 中间激活和 backward 状态}；LLM 推理更怕\key{长上下文 prefill 的计算量}和\key{decode 阶段不断增长的 KV cache}。
\end{intubox}

\subsection{六、为什么说 FlashAttention 是“实现优化”，不是“模型改造”？}

这是第一节最重要的概念边界。

标准 attention 的数学目标是：

\[
O_i = \sum_j \mathrm{softmax}\left(\frac{q_i k_j^\top}{\sqrt{d}} + M_{ij}\right)v_j
\]

FlashAttention 不改变这个公式。它不是说“只看一部分 token”，也不是说“用近似 softmax”。它做的是：

\begin{itemize}[leftmargin=1.5em]
  \item 把 $Q,K,V$ 分块搬到更快的小内存里计算；
  \item 用 online softmax 在分块扫描时维持数值稳定；
  \item 尽量不把完整 $n \times n$ score/probability matrix 写回 HBM；
  \item backward 阶段通过重算部分中间量，减少 forward 需要保存的状态。
\end{itemize}

\begin{warnbox}
\textbf{所以你看别人代码时要分清两类东西：}\\
\textbf{FlashAttention / SDPA backend} 通常是在\key{更高效地计算同一个 dense/causal attention}；\\
\textbf{sliding window / sparse / block mask} 则是在\key{改变哪些 token pair 允许互相注意}。前者主要是实现优化，后者会改变 attention pattern。
\end{warnbox}

\subsection{七、看代码时应该抓哪些关键词？}

第一节虽然不深入接口，但你现在可以先建立代码阅读雷达。看到下面这些词时，你应该意识到它们和 attention backend 或瓶颈有关：

\begin{center}
{\small
\begin{tabular}{p{4.2cm}p{8.6cm}}
\hline
\textbf{关键词} & \textbf{你应该想到什么} \\
\hline
\texttt{scaled\_dot\_product\_attention} & PyTorch SDPA 统一入口，底层可能选择 math、memory-efficient 或 flash 类 backend。 \\
\texttt{attn\_implementation} & Hugging Face 常见配置，例如 \texttt{eager}、\texttt{sdpa}、\texttt{flash\_attention\_2}。 \\
\texttt{flash\_attn\_func} & 直接调用 \texttt{flash-attn} 包的 dense attention kernel。 \\
\texttt{flash\_attn\_varlen\_func} & 变长序列 attention，通常会配合 \texttt{cu\_seqlens} 和 \texttt{max\_seqlen}。 \\
\texttt{is\_causal} / \texttt{causal} & 是否使用因果 mask，即 token 不能看未来。 \\
\texttt{attention\_mask} & 可能是 padding mask、causal mask 或业务自定义 mask；复杂 mask 会影响是否能命中高性能 backend。 \\
\texttt{sliding\_window} & 局部窗口 attention，常见于长上下文或视频 token 场景。 \\
\texttt{head\_dim} & 某些 kernel 对 head dimension 有限制，太大或不对齐可能 fallback。 \\
\hline
\end{tabular}
}
\end{center}

\begin{summarybox}
\textbf{代码阅读第一原则：} 不要只看函数名里有没有 “flash”。你要同时看 \key{Q/K/V shape、dtype、mask 类型、causal 标志、head dim、是否变长序列}。这些条件共同决定它到底能不能走高性能路径。
\end{summarybox}

\subsection{八、这一节的最终账本}

把这一节压缩成一张表：

\begin{center}
{\small
\begin{tabular}{p{3.2cm}p{4.2cm}p{5.3cm}}
\hline
\textbf{维度} & \textbf{标准 attention 的成本} & \textbf{工程含义} \\
\hline
计算 & $QK^\top$ 和 $PV$ 都是 $O(n^2d)$ & 长序列和视频 token 会显著放大计算量。 \\
显存 & score/probability 可能是 $B \times H \times n \times n$ & 朴素训练实现会被 $n^2$ 中间矩阵拖垮。 \\
IO & $S/P$ 大矩阵在 HBM 与片上高速内存之间反复搬运 & 速度可能受 memory bandwidth 限制；FlashAttention 重点减少 HBM 读写，而不是减少 attention 的数学语义。 \\
训练 & 要保存或重算 backward 所需中间状态 & FlashAttention、checkpointing、mixed precision 常组合使用。 \\
推理 & prefill 受长 prompt 影响，decode 受 KV cache 影响 & FlashAttention 主要帮助 prefill / dense attention；KV cache 另有账本。 \\
mask & 不同 mask 改变 token pair 可见性 & 简单 dense/causal 适合 SDPA/FlashAttention，复杂 mask 后面看 FlexAttention。 \\
\hline
\end{tabular}
}
\end{center}

\begin{conceptbox}
\textbf{Stage 0 最终结论：} 标准 attention 的问题不是公式难，而是工程代价大。$n^2$ 同时体现在\key{计算量、显存中间矩阵和 HBM IO} 上。FlashAttention 的意义，就是在不改变 dense/causal attention 数学结果的前提下，减少 $n \times n$ 中间矩阵物化和内存读写；online softmax 则解释了为什么分块扫描 score tile 以后，仍然能得到标准 softmax 的结果。你先把这个问题定义吃透，后面学 tiling、SDPA、FlexAttention 才不会散。
\end{conceptbox}

\subsection{九、你学完这一节后应该能立刻回答的问题}

\begin{itemize}[leftmargin=1.5em]
  \item 标准 attention 里 $QK^\top$ 和 $PV$ 的形状分别是什么？
  \item 为什么 attention 的 token pair 数是 $O(n^2)$？
  \item 为什么只说 FLOPs 不够，还要看 HBM / SRAM IO？
  \item 为什么 8k 到 32k sequence length 会让 attention matrix 变成 16 倍？
  \item online softmax 里 running max 和 running sum 分别在维护什么？
  \item 如果新 block 的最大值更大，为什么旧分母和旧输出累加器都要缩放？
  \item 训练 OOM 时，attention 相关的中间状态主要有哪些？
  \item LLM 推理里的 prefill 和 decode 为什么不是同一种瓶颈？
  \item FlashAttention 为什么是 exact attention 的实现优化，而不是 sparse/approximate attention？
\end{itemize}

% ============================================================
\subsection{参考资料}
% ============================================================

\begingroup
\small
\sloppy
\begin{itemize}[leftmargin=1.5em]
  \item Vaswani 等，\emph{Attention Is All You Need}, arXiv:1706.03762。原始 Transformer 论文，给出 scaled dot-product attention 和 multi-head attention 的基本形式。\\
  \url{https://arxiv.org/abs/1706.03762}
  \item Dao 等，\emph{FlashAttention: Fast and Memory-Efficient Exact Attention with IO-Awareness}, arXiv:2205.14135。本文明确指出标准 attention 的瓶颈不仅是 FLOPs，也包括 HBM 读写和 $n^2$ 中间矩阵物化；同时使用 tiling 与 online softmax 来避免显式保存完整 attention matrix。\\
  \url{https://arxiv.org/abs/2205.14135}
  \item Milakov 和 Gimelshein，\emph{Online normalizer calculation for softmax}, arXiv:1805.02867。介绍 online normalizer / online softmax 的递推思想，是理解分块 softmax 归一化的重要背景资料。\\
  \url{https://arxiv.org/abs/1805.02867}
  \item Dao，\emph{FlashAttention-2: Faster Attention with Better Parallelism and Work Partitioning}, arXiv:2307.08691。作为后续版本，用来理解 FlashAttention 从省显存进一步走向更高 GPU 利用率。\\
  \url{https://arxiv.org/abs/2307.08691}
  \item PyTorch \texttt{scaled\_dot\_product\_attention} 文档：PyTorch 2.x 的 SDPA 统一接口，以及 math、memory-efficient、FlashAttention 等 backend 语义。\\
  \url{https://docs.pytorch.org/docs/stable/generated/torch.nn.functional.scaled_dot_product_attention.html}
  \item PyTorch CUDA notes：Scaled Dot Product Attention backend 控制与 CUDA 相关说明。\\
  \url{https://docs.pytorch.org/docs/stable/notes/cuda.html#scaled-dot-product-attention}
  \item PyTorch \texttt{MultiheadAttention} 文档：标准多头 attention 模块和 PyTorch 内部对 SDPA 快路径的说明。\\
  \url{https://pytorch.org/docs/stable/generated/torch.nn.MultiheadAttention.html}
  \item NVIDIA CUDA C++ Programming Guide：CUDA memory hierarchy、global memory、shared memory、registers 等 GPU 内存层级说明。\\
  \url{https://docs.nvidia.com/cuda/cuda-c-programming-guide/index.html#memory-hierarchy}
  \item NVIDIA CUDA C++ Best Practices Guide：memory optimizations 章节，用于理解内存访问、带宽和数据复用对 kernel 性能的影响。\\
  \url{https://docs.nvidia.com/cuda/cuda-c-best-practices-guide/index.html#memory-optimizations}
  \item NVIDIA GPU Performance Background User's Guide：GPU 性能背景、内存带宽和计算吞吐之间关系的官方说明。\\
  \url{https://docs.nvidia.com/deeplearning/performance/dl-performance-gpu-background/index.html}
\end{itemize}
\endgroup

% ============================================================
% Stage 1：掌握 FlashAttention 的核心直觉，不钻 kernel
% ============================================================

\section{Stage 1：掌握 FlashAttention 的核心直觉，不钻 kernel}

\begin{conceptbox}
\textbf{最终口径：} FlashAttention 本质上是为了提升标准 attention 的执行效率并降低显存占用。它的核心逻辑是\key{分块处理 $Q/K/V$，让局部 score tile 在 GPU 片上高速内存中完成 scale、mask、online softmax 和对 $V$ 的累加，避免把完整的 $S$ 和 $P$ 这两个 $n \times n$ 中间矩阵写入 HBM。}
\end{conceptbox}

\begin{conceptbox}
\textbf{这一节的目标：} 你不需要会写 FlashAttention kernel，但要能建立正确直觉。学完这一节，你至少要能回答五件事：\\
1. \key{FlashAttention 为什么是 exact attention，而不是近似 attention}；\\
2. \key{它为什么能省显存：少物化 $S/P$ 这两个 $n \times n$ 中间矩阵}；\\
3. \key{它为什么可能更快：减少 HBM 读写，让更多计算在片上小工作区完成}；\\
4. \key{tiling 和 online softmax 在里面分别解决什么问题}；\\
5. \key{作为 AIGC 工程使用者，什么时候应该优先想到它，什么时候不要神化它。}
\end{conceptbox}

\note{这一节只讲“看懂和使用所需的核心直觉”，不讲 CUDA/Triton kernel 源码、不讲 warp/block 调度、不推完整 IO complexity 证明。}

\subsection{一、先给结论：FlashAttention 到底是什么？}

\begin{conceptbox}
\textbf{一句话版：} FlashAttention 是一种 \key{IO-aware 的 exact attention 实现}。它不改变 attention 公式，不少看 token，不近似 softmax；它只是把标准 attention 的计算顺序改成更适合 GPU 内存层级的方式，避免把完整 $n \times n$ attention score/probability matrix 写回 HBM。
\end{conceptbox}

标准 attention 的数学目标仍然是：

\[
O=\mathrm{softmax}\left(\frac{QK^\top}{\sqrt{d}}+M\right)V
\]

FlashAttention 算的还是这个东西。它没有把 dense attention 变成 sparse attention，也没有把 causal mask 改成别的 mask。它主要改变的是\key{实现路径}：

\begin{center}
{\small
\begin{tabular}{p{4.1cm}p{8.6cm}}
\hline
\textbf{朴素 attention} & 先算完整 $S=QK^\top$，写回 HBM；再读 $S$ 做 softmax 得到 $P$，写回 HBM；再读 $P,V$ 算 $O=PV$。 \\
\hline
\textbf{FlashAttention} & 分块读取 $Q,K,V$，在片上小工作区里算当前 tile 的 score；用 online softmax 合并局部结果；不保存完整 $S/P$，最后只写回 $O$。 \\
\hline
\end{tabular}
}
\end{center}

\begin{intubox}
\textbf{你可以这样记：} FlashAttention 不是“少算 attention”，而是“别把巨大的中间表写来写去”。它省的是 $S/P$ 的物化和 HBM 往返，不是改变 attention 的语义。
\end{intubox}

\subsection{二、FlashAttention 想解决的不是公式，而是执行账本}

Stage 0 里已经建立过账本：朴素实现里的大问题不是 $Q,K,V,O$ 本身，而是中间矩阵：

\[
S=QK^\top \in \mathbb{R}^{n\times n}
\]

\[
P=\mathrm{softmax}(S) \in \mathbb{R}^{n\times n}
\]

如果序列长度 $n$ 很大，$S$ 和 $P$ 就会非常大。更麻烦的是，它们不是只“存在一下”，而是可能经历多次 HBM 读写：

\begin{center}
\begin{tabular}{p{13.2cm}}
\hline
读 $Q,K$，算完整 $S$，写回 HBM \\
$\downarrow$ \\
读 $S$，做 mask / softmax / dropout，写完整 $P$ 回 HBM \\
$\downarrow$ \\
读 $P,V$，算 $O=PV$，写 $O$ 回 HBM \\
\hline
\end{tabular}
\end{center}

\begin{warnbox}
\textbf{这里的关键不是“矩阵乘很慢”这么简单。} GPU 的计算吞吐很强，但 HBM 和片上内存之间搬大矩阵仍然贵。FlashAttention 论文强调的 IO-aware，就是算法设计时显式关心 HBM 读写量，而不是只看 FLOPs。
\end{warnbox}

\subsection{三、核心直觉 1：tiling，把大表切成小块来算}

FlashAttention 的第一个直觉是 \key{tiling}，也就是分块。

不要一次性构造完整的：

\[
S\in \mathbb{R}^{n\times n}
\]

而是把 $Q$ 按 query block 切，把 $K,V$ 按 key/value block 切。每次只拿一小块进入片上高速工作区：

\[
Q_i \in \mathbb{R}^{B_q\times d},\quad K_j,V_j \in \mathbb{R}^{B_k\times d}
\]

然后算当前局部 score tile：

\[
S_{ij}=Q_iK_j^\top \in \mathbb{R}^{B_q\times B_k}
\]

这个 tile 很小，可以在 shared memory / registers / cache 等片上层级里高效处理。处理完之后，不把它作为完整 attention matrix 的一部分写回 HBM，而是马上参与 softmax 和 $V$ 的累加。

\begin{center}
{\small
\begin{tabular}{p{3.2cm}p{4.4cm}p{5.3cm}}
\hline
\textbf{对象} & \textbf{朴素理解} & \textbf{FlashAttention 直觉} \\
\hline
$S$ & 一个完整 $n\times n$ 大矩阵 & 只临时出现一个小 tile，算完就合并。 \\
$P$ & 一个完整 softmax 概率矩阵 & 通常不完整物化，只用局部概率贡献更新输出。 \\
$O$ & 最后输出 & 对每个 query block 维护累加结果，最后写回 HBM。 \\
\hline
\end{tabular}
}
\end{center}

\begin{intubox}
\textbf{类比：} 朴素 attention 像是先把整张巨大的表格填完、存档、再拿出来统计；FlashAttention 像是分批读数据，每批读完立刻更新统计量，最终得到同一个统计结果。
\end{intubox}

\subsection{四、核心直觉 2：online softmax，让分块结果仍然等价}

\begin{conceptbox}
\textbf{先说人话：} online softmax 解决的是一个很具体的问题：\key{FlashAttention 分块看 score，但 softmax 又需要整行 score 的全局最大值和全局分母。} 所以它不能只对每个小块各自 softmax，否则就变成“局部 softmax”，语义会错；它必须边看 block，边维护足够的统计量，最后得到和一次性看完整行一样的 softmax。
\end{conceptbox}

普通稳定 softmax 要先看完整一行 logits。比如一行 score 是：

\[
[1,2,5]
\]

稳定 softmax 会先找全局最大值：

\[
m=5
\]

然后再算：

\[
\frac{e^{1-5},\ e^{2-5},\ e^{5-5}}{e^{1-5}+e^{2-5}+e^{5-5}}
\]

问题是 FlashAttention 不想先把完整 $S=QK^\top$ 算出来并存下来。它可能先看到：

\[
[1,2]
\]

后面才看到：

\[
[5]
\]

这时最大的困难是：\key{前面只看到 $[1,2]$ 的时候，你并不知道后面会不会出现更大的 $5$。}

\subsubsection*{1. online softmax 只记三类“摘要”，不记完整 score}

online softmax 不保存完整 score 行，而是维护下面三个东西：

\begin{center}
{\small
\begin{tabular}{p{2.0cm}p{4.7cm}p{5.8cm}}
\hline
\textbf{记号} & \textbf{人话含义} & \textbf{为什么需要} \\
\hline
$m$ & 到目前为止见过的最大 score & softmax 为了数值稳定，要统一减去最大值。 \\
$l$ & 到目前为止的 softmax 分母 & 最后要用它做归一化。 \\
$acc$ 或 $o$ & 到目前为止的 value 加权和 & attention 最终要的是 softmax 权重乘 $V$ 的结果。 \\
\hline
\end{tabular}
}
\end{center}

\begin{intubox}
\textbf{通俗类比：} 不把所有明细账单都存下来，只随身带三个摘要：目前最大一笔、目前总权重、目前加权结果。新账单来一批，就更新这三个摘要。
\end{intubox}

\subsubsection*{2. 为什么后面来了更大的 score，旧结果要“缩小”？}

先看第一块 $[1,2]$。当时最大值是：

\[
m_{old}=2
\]

所以旧分母是按 $m=2$ 算的：

\[
l_{old}=e^{1-2}+e^{2-2}
\]

现在第二块来了一个更大的 score：

\[
5
\]

新的最大值变成：

\[
m_{new}=5
\]

正确的全局分母应该按 $m=5$ 重新表达：

\[
e^{1-5}+e^{2-5}+e^{5-5}
\]

注意前两项可以从旧结果变过来：

\[
e^{1-5}=e^{1-2}\cdot e^{2-5}
\]

\[
e^{2-5}=e^{2-2}\cdot e^{2-5}
\]

也就是说，旧分母整体乘上：

\[
e^{m_{old}-m_{new}}=e^{2-5}
\]

\begin{summarybox}
\textbf{最直白的理解：} 后面来了更大的 score，softmax 的参照物从 $2$ 变成 $5$。以前那些 score 相对就“没那么大”了，所以旧分母、旧加权和都要整体缩小，再把新 block 加进去。
\end{summarybox}

\subsubsection*{3. 公式只是把这件事写紧凑一点}

对一行 logits $x=[x_1,x_2,\dots,x_n]$，稳定 softmax 是：

\[
m=\max_j x_j
\]

\[
\mathrm{softmax}(x_i)=\frac{e^{x_i-m}}{\sum_j e^{x_j-m}}
\]

如果新 block 的局部最大值是 $m_{block}$，局部分母是 $l_{block}$，则合并时：

\[
m_{new}=\max(m_{old},m_{block})
\]

\[
l_{new}
=
e^{m_{old}-m_{new}}l_{old}
+
e^{m_{block}-m_{new}}l_{block}
\]

attention 还需要输出加权和，所以也维护：

\[
o=\sum_j e^{x_j-m}v_j
\]

新 block 合并时，旧输出也要按同样逻辑缩放：

\[
o_{new}
=
e^{m_{old}-m_{new}}o_{old}
+
e^{m_{block}-m_{new}}o_{block}
\]

最后再归一化：

\[
O=\frac{o}{l}
\]

\begin{summarybox}
\textbf{online softmax 在 FlashAttention 里的意义：} 它让 FlashAttention 可以\key{一边分块扫描 score，一边保持全局 softmax 语义}。没有 online softmax，tiling 只能得到局部 softmax；有了 online softmax，分块计算仍然等价于标准 softmax attention。
\end{summarybox}

\subsection{五、核心直觉 3：fusion，别让中间结果落地}

FlashAttention 另一个很重要的工程直觉是 \key{fusion}：把原本分散的步骤尽量融合在一个 kernel 或少数 kernel 的执行流程里。

朴素实现可以理解成多个阶段：

\begin{center}
\begin{tabular}{p{13.2cm}}
\hline
matmul 产生 $S$ $\rightarrow$ mask / scale $\rightarrow$ softmax $\rightarrow$ dropout $\rightarrow$ matmul 产生 $O$ \\
\hline
\end{tabular}
\end{center}

如果每个阶段都把中间结果写回 HBM，再被下一个阶段读出来，HBM IO 就会很重。FlashAttention 的思路是尽量在处理 tile 时完成这些操作：

\begin{itemize}[leftmargin=1.5em]
  \item scale 和 mask 直接作用在当前 score tile 上；
  \item softmax 统计量用 online 方式更新；
  \item 当前 tile 对 $V$ 的贡献立刻累加到输出 accumulator；
  \item 当前 $S_{ij}$ 和局部 $P_{ij}$ 用完就丢，不写成完整大矩阵。
\end{itemize}

\begin{intubox}
\textbf{使用者视角的理解：} fusion 不是为了改变模型，而是为了减少“中间 tensor 落地”。你在代码里可能只看到一个 attention 调用，但底层可能已经把 scale、mask、softmax、dropout、$PV$ 融合成更高效的路径。
\end{intubox}

\subsection{六、核心直觉 4：backward 也要省，必要时重算}

\begin{conceptbox}
\textbf{先说人话：} backward recompute 解决的是训练显存问题。朴素 attention 为了反向传播方便，会在 forward 时保存完整 softmax probability 矩阵 $P$；FlashAttention 觉得这个 $P$ 太大，于是选择\key{forward 不保存完整 $P$，backward 需要哪一小块就重新算哪一小块。}
\end{conceptbox}

训练时，forward 算完还不够，backward 要根据中间结果算梯度。朴素 attention 常见思路是：

\begin{center}
\begin{tabular}{p{13.2cm}}
\hline
\texttt{forward}: 算 $S=QK^\top$，算 $P=\mathrm{softmax}(S)$，把完整 $P$ 保存下来。 \\
$\downarrow$ \\
\texttt{backward}: 直接读取保存好的 $P$，用它计算 $dQ,dK,dV$。 \\
\hline
\end{tabular}
\end{center}

这样做的好处是 backward 方便，坏处是 $P$ 很大：

\[
P\in \mathbb{R}^{B\times H\times n\times n}
\]

长序列训练时，保存完整 $P$ 就像把一整本巨大的“中间过程复印件”都留下来，非常占显存。

\subsubsection*{1. FlashAttention 的选择：不保存完整复印件}

FlashAttention forward 的态度是：

\begin{center}
\begin{tabular}{p{13.2cm}}
\hline
我不保存完整 $S$，也不保存完整 $P$；我只保存 backward 必须用到的小摘要，例如输出 $O$ 和 softmax normalizer 相关信息。 \\
\hline
\end{tabular}
\end{center}

等 backward 真需要某个局部 probability 时，再分块重新算：

\begin{center}
\begin{tabular}{p{13.2cm}}
\hline
\texttt{backward}: 重新读取对应的 $Q$ block、$K$ block、$V$ block，重算局部 $S_{ij}$，恢复局部 softmax contribution，用完就丢。 \\
\hline
\end{tabular}
\end{center}

\begin{intubox}
\textbf{生活类比：} 朴素 attention 像是第一次算账时，把每一页明细都复印保存，之后查账直接翻复印件；FlashAttention 像是只保存总额、关键摘要和原始账本位置，之后需要哪一页细节，就回原始账本重新算那一页。
\end{intubox}

\subsubsection*{2. backward 不是完整粗暴地重跑 forward}

这里要避免一个误解：FlashAttention backward 不是简单地“完整重跑一遍 forward”。更准确地说，它是\key{按照 tile 重算局部中间量}。

它会在 backward 中局部重算：

\begin{itemize}[leftmargin=1.5em]
  \item 当前 tile 的 $QK^\top$ score；
  \item 当前 tile 对应的 softmax probability contribution；
  \item 当前 tile 对 $dQ,dK,dV$ 的梯度贡献。
\end{itemize}

但是它仍然不会把完整 $S$ 或完整 $P$ 保存下来。它的模式是：

\begin{center}
\begin{tabular}{p{13.2cm}}
\hline
需要哪一小块 $\rightarrow$ 重算哪一小块 $\rightarrow$ 用这块算梯度贡献 $\rightarrow$ 用完丢掉 $\rightarrow$ 处理下一小块。 \\
\hline
\end{tabular}
\end{center}

\subsubsection*{3. 为什么重算是划算的？}

这就是一个 trade-off：

\begin{center}
{\small
\begin{tabular}{p{3.4cm}p{4.4cm}p{5.0cm}}
\hline
\textbf{策略} & \textbf{省什么} & \textbf{付出什么} \\
\hline
保存完整 $P$ & backward 少重算 & forward 激活显存大，$n^2$ 压力重。 \\
不保存完整 $P$，backward 分块重算局部量 & 显著减少 forward 需要保存的中间矩阵 & backward 多一些计算，但通常整体更划算。 \\
\hline
\end{tabular}
}
\end{center}

\begin{warnbox}
\textbf{不要把“重算”理解成坏事。} 在 GPU 上，很多时候\key{多算一小块}比\key{保存巨大中间矩阵、再反复从 HBM 读回来}更划算。这和 activation checkpointing 的系统思想很像：用可控的额外计算换显存。
\end{warnbox}

\begin{summarybox}
\textbf{最直白的总结：} FlashAttention backward 的核心不是神秘技巧，而是\key{forward 不存完整 attention probability，backward 需要时再分块重算局部 probability，用一点额外计算换掉巨大的显存占用。}
\end{summarybox}

\subsection{七、FlashAttention 为什么 exact？}

这是你看别人代码和面试解释时最容易被问到的点。

\begin{conceptbox}
\textbf{FlashAttention exact 的意思：} 在同样的 $Q,K,V$、同样的 mask、同样的 scale、同样的 dropout 设置下，它目标上计算的是同一个 attention 公式；差别只在实现顺序和浮点舍入误差。它不是 Linformer、Performer、local attention 或 block-sparse attention 这类改变 attention pattern 或近似 attention matrix 的方法。
\end{conceptbox}

可以用这张表来区分：

\begin{center}
{\small
\begin{tabular}{p{3.3cm}p{4.4cm}p{5.2cm}}
\hline
\textbf{类型} & \textbf{是否改变 token pair 可见性} & \textbf{你应该怎么理解} \\
\hline
FlashAttention & 不改变 & 更高效地算同一个 dense/causal attention。 \\
Sliding window attention & 改变 & 每个 token 只看局部窗口，attention pattern 变了。 \\
Block-sparse attention & 改变 & 只计算某些 block，通常不是完整 dense attention。 \\
Approximate kernel attention & 通常改变或近似 & 用近似技巧降低复杂度，数学结果不再严格等同普通 softmax attention。 \\
\hline
\end{tabular}
}
\end{center}

\begin{summarybox}
\textbf{一句话背下来：} FlashAttention 是 exact attention 的高性能实现；sparse / sliding window / approximate attention 是改变或近似 attention pattern 的算法。不要把它们混成一类。
\end{summarybox}

\subsection{八、FlashAttention 为什么省显存？为什么可能更快？}

从使用者角度，可以把收益拆成两类。

\subsubsection*{1. 省显存：少保存 $n^2$ 中间矩阵}

朴素 attention 最直观的显存压力来自：

\[
S,P \in \mathbb{R}^{B\times H\times n\times n}
\]

FlashAttention 避免完整物化这些矩阵，训练时 forward 保存的中间状态也更少，所以长序列训练或大 batch 下更不容易因为 attention 激活 OOM。

\begin{intubox}
\textbf{工程口径：} 如果你打开 FlashAttention 后发现显存下降，通常不是因为 $Q,K,V,O$ 消失了，而是因为完整 $S/P$ 这类大中间矩阵不再被保存或反复落地。
\end{intubox}

\subsubsection*{2. 可能更快：少搬数据，提高实际吞吐}

FlashAttention 不是简单地把 FLOPs 从 $O(n^2d)$ 变成 $O(n d)$。它仍然在算 dense attention，所以主要矩阵乘的计算量级仍然和 token pair 有关。

它快的原因更偏系统层面：

\begin{itemize}[leftmargin=1.5em]
  \item 减少 $S/P$ 在 HBM 中的写入和读取；
  \item 提高 $K,V$ tile 的复用；
  \item 把 scale、mask、softmax、dropout、$PV$ 等步骤融合，减少 kernel 往返和中间 tensor；
  \item 让更多中间计算停留在片上高速内存和寄存器里。
\end{itemize}

\begin{warnbox}
\textbf{不要误解：} FlashAttention 不会让 dense attention 的理论 token pair 数从 $n^2$ 消失。长序列仍然贵；它只是让标准 dense/causal attention 的实现更接近硬件友好的路径。
\end{warnbox}

\subsection{九、FlashAttention-1 / 2 / 3 作为使用者怎么理解？}

作为工程使用者，不需要深入版本论文的所有 kernel 细节，但需要知道版本演进大致在优化什么：

\begin{center}
{\small
\begin{tabular}{p{2.8cm}p{5.3cm}p{4.8cm}}
\hline
\textbf{版本} & \textbf{主要直觉} & \textbf{使用者应该关心什么} \\
\hline
FlashAttention & IO-aware exact attention，通过 tiling 和 online softmax 减少 HBM 读写和 $n^2$ 中间矩阵物化。 & 是否能启用、是否支持当前 dtype / head dim / mask。 \\
FlashAttention-2 & 更好的并行划分和 work partitioning，减少非 matmul 工作比例，提高 GPU 利用率。 & 新项目通常优先看到 \texttt{flash\_attention\_2}；关注版本兼容和性能差异。 \\
FlashAttention-3 & 面向 Hopper 等新硬件进一步利用异步和低精度能力。 & 关注硬件架构、CUDA/PyTorch/包版本是否支持，而不是手写 kernel。 \\
\hline
\end{tabular}
}
\end{center}

\begin{intubox}
\textbf{实际项目里怎么想：} 大多数时候你不需要手动选择“论文里的第几版算法”，而是通过 PyTorch SDPA、Hugging Face \texttt{attn\_implementation} 或 \texttt{flash-attn} 包间接使用。你真正要排查的是：\key{有没有命中高性能 backend、mask 是否支持、dtype/head dim 是否满足、版本是否兼容。}
\end{intubox}

\subsection{十、在 AIGC 场景里，它通常帮你解决什么？}

\begin{center}
{\small
\begin{tabular}{p{3.2cm}p{4.7cm}p{4.8cm}}
\hline
\textbf{场景} & \textbf{为什么 attention 贵} & \textbf{FlashAttention 的帮助} \\
\hline
LLM 长上下文训练 / prefill & prompt 长，dense causal attention 的 $n^2$ 压力明显。 & 减少 attention 中间矩阵显存和 HBM IO，常用于长上下文训练和 prefill 加速。 \\
DiT / 图像生成 & 图像 patch token 随分辨率增长。 & 对 dense self-attention 层通常有帮助，尤其是高分辨率或大 batch。 \\
视频生成 & 帧数 $\times$ 空间 patch token，token 数更容易爆炸。 & 如果仍是 dense/causal attention，可减少中间矩阵压力；但视频常还需要窗口或稀疏 pattern。 \\
多图多文本条件 & token 来源复杂，mask 可能更复杂。 & 普通 dense/causal 部分适合；复杂结构化 mask 可能需要后续看 FlexAttention。 \\
\hline
\end{tabular}
}
\end{center}

\begin{summarybox}
\textbf{使用者判断原则：} 如果你的 attention 仍然是标准 dense 或 causal，优先想到 SDPA / FlashAttention；如果你的主要需求是“哪些 token 能看哪些 token”的复杂规则，那就不只是 FlashAttention 问题，后面要看 FlexAttention 或其他 sparse/block mask 方案。
\end{summarybox}

\subsection{十一、什么时候不要神化 FlashAttention？}

FlashAttention 很重要，但不是所有慢和 OOM 都能靠它解决。

\begin{itemize}[leftmargin=1.5em]
  \item \textbf{如果瓶颈是 MLP、卷积、VAE、数据加载或通信：} 开 FlashAttention 不会直接解决这些瓶颈。
  \item \textbf{如果 decode 阶段主要受 KV cache 容量或读取影响：} FlashAttention 可能帮一部分 attention 计算，但 KV cache 账本仍然要单独处理。
  \item \textbf{如果 mask 很复杂：} 可能无法命中普通 flash backend，或者需要 FlexAttention / block-sparse 方案。
  \item \textbf{如果 dtype、head dim、GPU 架构、包版本不满足要求：} 代码可能 fallback 到 math/eager 路径，实际速度和显存收益会变小。
  \item \textbf{如果序列很短：} attention 中间矩阵本来不大，FlashAttention 的收益可能没有长序列明显。
\end{itemize}

\begin{warnbox}
\textbf{看性能时不要只看开关。} 配置里写了 \texttt{flash\_attention\_2} 不代表一定命中最快路径；你还要看实际 backend、输入 shape、mask、dtype、head dim、训练/推理阶段和硬件。
\end{warnbox}

\subsection{十二、看别人代码时应该形成的 mental model}

看到 attention 代码时，你可以按下面顺序理解：

\begin{enumerate}[leftmargin=1.7em]
  \item \textbf{语义层：} 这是 dense attention、causal attention、sliding window，还是复杂 mask？
  \item \textbf{实现层：} 它走的是 eager/math、PyTorch SDPA、\texttt{flash-attn}，还是框架自定义 kernel？
  \item \textbf{输入层：} $Q,K,V$ 的 shape 是什么？batch、heads、sequence length、head dim 分别是多少？
  \item \textbf{约束层：} dtype、GPU、dropout、mask、head dim 是否满足高性能 backend 条件？
  \item \textbf{收益层：} 当前场景是训练、prefill 还是 decode？瓶颈真的是 attention 中间矩阵和 HBM IO 吗？
\end{enumerate}

\begin{conceptbox}
\textbf{代码阅读一句话：} FlashAttention 相关代码不要只看函数名，要同时看 \key{attention 语义、backend、shape、mask、dtype、head dim 和运行阶段}。这些条件共同决定它到底是在高效地算标准 attention，还是已经 fallback / 改变了 attention pattern。
\end{conceptbox}

\subsection{十三、这一节的最终账本}

\begin{center}
{\small
\begin{tabular}{p{3.0cm}p{4.7cm}p{5.0cm}}
\hline
\textbf{关键词} & \textbf{核心含义} & \textbf{你要记住的工程结论} \\
\hline
Exact attention & 数学目标仍然是标准 softmax attention。 & 不是近似、不是稀疏、不是模型结构变化。 \\
Tiling & 把 $Q,K,V$ 分块，局部处理 score tile。 & 不构造完整 $S$，减少大中间矩阵落地。 \\
Online softmax & 边扫描 block，边维护 running max / sum / output accumulator。 & 分块后仍然能得到全局 softmax 结果。 \\
Fusion & scale、mask、softmax、dropout、$PV$ 尽量融合处理。 & 少写少读中间 tensor，降低 HBM IO。 \\
Backward recompute & 不保存完整 $P$，反向分块重算局部量。 & 用一些计算换显存，训练更友好。 \\
限制条件 & dtype、mask、head dim、GPU、版本都会影响 backend。 & 开了配置不代表一定命中 flash 路径。 \\
\hline
\end{tabular}
}
\end{center}

\begin{summarybox}
\textbf{Stage 1 最终结论：} FlashAttention 的核心不是“一个神秘 kernel”，而是一个很清楚的系统优化思路：\key{把标准 attention 的大矩阵计算改成分块扫描，用 online softmax 保持 exact，用 fusion 和重算减少 $n^2$ 中间矩阵物化和 HBM IO。} 你不需要钻 CUDA kernel，也应该能解释它为什么 exact、为什么省显存、为什么常常更快，以及为什么复杂 mask 场景还需要后续学习 FlexAttention。
\end{summarybox}

\subsection{十四、你学完这一节后应该能立刻回答的问题}

\begin{itemize}[leftmargin=1.5em]
  \item FlashAttention 为什么是 exact attention，而不是 approximate attention？
  \item tiling 解决的是什么问题？为什么它不是简单地把完整 $S$ 放进 SRAM？
  \item online softmax 为什么能让分块 softmax 等价于全局 softmax？
  \item FlashAttention 为什么能减少显存峰值？它主要少保存了什么？
  \item FlashAttention 为什么可能更快？它减少的是哪类 IO？
  \item backward 里为什么可以接受重算一部分中间量？
  \item 为什么开了 \texttt{flash\_attention\_2} 不代表一定命中最快路径？
  \item 普通 dense/causal attention 和复杂 mask attention 的优化思路有什么区别？
\end{itemize}

% ============================================================
\subsection{参考资料}
% ============================================================

\begingroup
\small
\sloppy
\begin{itemize}[leftmargin=1.5em]
  \item Dao 等，\emph{FlashAttention: Fast and Memory-Efficient Exact Attention with IO-Awareness}, arXiv:2205.14135。原始 FlashAttention 论文，提出 IO-aware exact attention，通过 tiling 和 online softmax 减少 HBM 访问与 $n^2$ 中间矩阵物化。\\
  \url{https://arxiv.org/abs/2205.14135}
  \item Dao，\emph{FlashAttention-2: Faster Attention with Better Parallelism and Work Partitioning}, arXiv:2307.08691。用于理解后续版本如何通过更好的并行划分和 work partitioning 提升 GPU 利用率。\\
  \url{https://arxiv.org/abs/2307.08691}
  \item Milakov 和 Gimelshein，\emph{Online normalizer calculation for softmax}, arXiv:1805.02867。online softmax / online normalizer 的背景论文，有助于理解分块 softmax 如何保持全局归一化。\\
  \url{https://arxiv.org/abs/1805.02867}
  \item Dao-AILab \texttt{flash-attention} 官方仓库：FlashAttention 官方实现、安装说明、接口、硬件与版本支持信息。\\
  \url{https://github.com/Dao-AILab/flash-attention}
  \item PyTorch \texttt{scaled\_dot\_product\_attention} 文档：PyTorch 2.x SDPA 统一入口以及 FlashAttention、memory-efficient、math backend 的使用语义。\\
  \url{https://docs.pytorch.org/docs/stable/generated/torch.nn.functional.scaled_dot_product_attention.html}
  \item PyTorch CUDA notes：SDPA backend 控制、CUDA attention backend 相关说明，用来理解实际代码中为什么可能 fallback。\\
  \url{https://docs.pytorch.org/docs/stable/notes/cuda.html#scaled-dot-product-attention}
  \item NVIDIA CUDA C++ Programming Guide：CUDA memory hierarchy，用于理解 HBM/global memory、shared memory、registers 等内存层级。\\
  \url{https://docs.nvidia.com/cuda/cuda-c-programming-guide/index.html#memory-hierarchy}
\end{itemize}
\endgroup

% ============================================================
% Stage 4：再学 FlexAttention，重点放在“什么时候需要”
% ============================================================

\section{Stage 4：再学 FlexAttention，重点放在“什么时候需要”}

\begin{conceptbox}
\textbf{最终口径：} FlexAttention 不是“比 FlashAttention 更高级的新默认”，而是当你的 attention 已经不是普通 dense / causal 模式，出现\key{自定义 token 可见性、自定义 pre-softmax score 修改、sliding window、prefix-LM、document mask、sample packing、block-sparse、多模态分区规则}时，用 \texttt{score\_mod} / \texttt{mask\_mod} 把规则写出来，再让 PyTorch 通过 \texttt{torch.compile} 尽量生成 fused、可利用稀疏性的高性能 attention kernel。
\end{conceptbox}

\note{这一节重点是“什么时候需要 FlexAttention”，不是教你手写 Triton/CUDA kernel，也不是把所有 attention 变体都实现一遍。你要先能判断：普通 SDPA / FlashAttention 够不够；如果不够，是不是该上 FlexAttention。}

\subsection{一、先给结论：什么时候需要 FlexAttention？}

如果一句话判断：

\begin{summarybox}
\textbf{普通 dense / causal attention：} 优先用 PyTorch SDPA / FlashAttention。\\
\textbf{复杂 attention pattern：} 如果你需要控制“哪些 token 能看哪些 token”，或者需要在 softmax 前对 score 做自定义修改，并且还想尽量保留 fused kernel 的性能，就开始考虑 FlexAttention。
\end{summarybox}

换句话说，FlexAttention 解决的不是“标准 attention 怎么更快”这个问题；这个问题主要由 SDPA / FlashAttention 负责。FlexAttention 解决的是：

\begin{center}
\begin{tabular}{p{13.2cm}}
\hline
我现在的 attention 规则太灵活，普通 flash backend 支持不了；但如果退回 eager / dense mask，又慢又吃显存。有没有办法用 PyTorch 写规则，同时让底层生成接近 fused attention 的实现？ \\
\hline
\end{tabular}
\end{center}

\begin{intubox}
\textbf{你可以这样记：} FlashAttention 主要优化“怎么算标准 attention”；FlexAttention 主要解决“我想定义不标准的 attention 规则，但又不想为了每一种规则手写 kernel”。
\end{intubox}

\subsection{二、先和 FlashAttention 分清边界}

\begin{center}
{\small
\begin{tabular}{p{3.0cm}p{5.0cm}p{5.0cm}}
\hline
\textbf{问题} & \textbf{优先想到 FlashAttention / SDPA} & \textbf{开始考虑 FlexAttention} \\
\hline
attention 语义 & 普通 dense self-attention 或 causal attention。 & token pair 可见性有自定义规则。 \\
mask 类型 & causal mask、简单 padding mask，且 backend 支持。 & sliding window、prefix-LM、document boundary、block-sparse、多规则组合。 \\
score 修改 & 没有额外 pre-softmax 修改，或框架已有内置支持。 & ALiBi / relative bias / tanh soft cap 等需要自定义 \texttt{score\_mod}。 \\
性能目标 & 更快地算标准公式，减少 $S/P$ 物化和 HBM IO。 & 自定义规则下仍希望 fused、少物化、能跳过 masked block。 \\
工程复杂度 & 配置或调用现成接口即可。 & 需要写 \texttt{score\_mod} / \texttt{mask\_mod}，并关注编译、缓存和 mask 稀疏结构。 \\
\hline
\end{tabular}
}
\end{center}

\begin{warnbox}
\textbf{不要把 FlexAttention 理解成“所有场景都更快”。} 如果你的 attention 本来就是普通 causal LLM attention，那么成熟的 SDPA / FlashAttention 往往是更直接、更稳的选择。FlexAttention 的优势出现在\key{规则复杂、普通 backend 不支持、但规则又有结构可利用}的时候。
\end{warnbox}

\subsection{三、FlexAttention 的最小心智模型}

标准 attention 可以写成：

\[
S_{ij}=\frac{q_i k_j^\top}{\sqrt d}
\]

\[
P_{ij}=\mathrm{softmax}(S_{ij})
\]

\[
O_i=\sum_j P_{ij}v_j
\]

FlexAttention 给你的关键能力是：在 score 进入 softmax 之前，插入可编程规则。

\subsubsection*{1. \texttt{score\_mod}：改 score 的数值}

\texttt{score\_mod} 的直觉是：对每一个 query-key pair 的 score，在 softmax 前做一次自定义变换。

概念上可以理解成：

\begin{lstlisting}[language=Python]
def score_mod(score, b, h, q_idx, kv_idx):
    return score
\end{lstlisting}

参数含义是：

\begin{center}
{\small
\begin{tabular}{p{2.3cm}p{10.2cm}}
\hline
\textbf{参数} & \textbf{含义} \\
\hline
\texttt{score} & 当前 query token 和 key token 的点积 score。 \\
\texttt{b} & 当前 batch index。 \\
\texttt{h} & 当前 head index。 \\
\texttt{q\_idx} & 当前 query 位置。 \\
\texttt{kv\_idx} & 当前 key/value 位置。 \\
\hline
\end{tabular}
}
\end{center}

它适合表达：

\begin{itemize}[leftmargin=1.5em]
  \item 给不同距离加 bias，例如相对位置 bias 或 ALiBi 风格 bias；
  \item 对 score 做 soft cap，例如某些模型里的 tanh soft-capping；
  \item 对不同 head、不同 token 区域使用不同的 pre-softmax 修改。
\end{itemize}

\begin{warnbox}
\textbf{注意：} 你也可以在 \texttt{score\_mod} 里把不允许的位置改成 $-\infty$，这样数学上可以实现 mask；但如果目标是利用 mask 的稀疏性跳过计算，通常应该用 \texttt{mask\_mod} + \texttt{BlockMask}。
\end{warnbox}

\subsubsection*{2. \texttt{mask\_mod}：定义哪些 token pair 可以参与计算}

\texttt{mask\_mod} 的直觉是：返回一个布尔值，告诉 kernel 当前 query-key pair 是否允许参与 attention。

概念上可以理解成：

\begin{lstlisting}[language=Python]
def mask_mod(b, h, q_idx, kv_idx):
    return q_idx >= kv_idx
\end{lstlisting}

其中 \texttt{True} 表示这个位置保留，\texttt{False} 表示这个位置被 mask 掉。

如果只是把 mask 写进 dense tensor，你仍然可能要处理一个很大的 $n\times n$ mask。FlexAttention 更重要的地方是：把 \texttt{mask\_mod} 交给 \texttt{create\_block\_mask}，生成 \texttt{BlockMask}，让底层知道哪些 block 可以整体跳过。

\subsubsection*{3. \texttt{BlockMask}：让 mask 的稀疏性真正变成性能收益}

\texttt{BlockMask} 可以理解成 block 级别的稀疏 mask 表示。它不是只告诉你“某个元素是否可见”，而是帮助 kernel 判断：

\begin{center}
\begin{tabular}{p{13.2cm}}
\hline
这一整块 query-key 区域是不是都不用算？如果整块都被 mask 掉，就可以跳过这块计算。 \\
\hline
\end{tabular}
\end{center}

\begin{intubox}
\textbf{最直白的理解：} \texttt{score\_mod} 更像“改分数”；\texttt{mask\_mod} 更像“规定谁能看谁”；\texttt{BlockMask} 更像“把谁能看谁这件事整理成适合 kernel 跳过整块计算的地图”。
\end{intubox}

\subsection{四、最典型的需要场景 1：sliding window / local attention}

长上下文里，如果每个 token 都看全部历史 token，attention 仍然有 $n^2$ token pair。很多模型会改成局部窗口：每个 token 只看附近一段上下文。

例如 causal sliding window：第 $i$ 个 query 只能看过去 $W$ 个 token：

\[
i-W \le j \le i
\]

用 FlexAttention 的 \texttt{mask\_mod} 可以表达成：

\begin{lstlisting}[language=Python]
from torch.nn.attention.flex_attention import flex_attention, create_block_mask

WINDOW = 1024

def sliding_causal_mask(b, h, q_idx, kv_idx):
    return (kv_idx <= q_idx) & (q_idx - kv_idx <= WINDOW)

block_mask = create_block_mask(
    sliding_causal_mask,
    B=None,
    H=None,
    Q_LEN=S,
    KV_LEN=S,
    device=query.device,
)

out = flex_attention(query, key, value, block_mask=block_mask)
\end{lstlisting}

\begin{summarybox}
\textbf{什么时候需要：} 如果你要的是标准 causal attention，就不需要 FlexAttention；如果你要的是\key{只看局部窗口、局部窗口再叠加 causal、局部窗口再叠加文档边界}，普通 flash backend 未必直接支持，这时 FlexAttention 就有价值。
\end{summarybox}

\subsection{五、最典型的需要场景 2：prefix-LM}

prefix-LM 的规则和普通 causal 不一样。常见形式是：

\begin{itemize}[leftmargin=1.5em]
  \item prefix 区域内部可以双向看；
  \item generation 区域可以看整个 prefix；
  \item generation 区域内部仍然 causal，只能看过去，不能看未来。
\end{itemize}

如果 prefix 长度是 $p$，可以粗略理解成：

\[
\text{allow}(i,j)=
\begin{cases}
\text{true}, & i < p \text{ and } j < p \\
\text{true}, & i \ge p \text{ and } j < p \\
\text{true}, & i \ge p \text{ and } p \le j \le i \\
\text{false}, & \text{otherwise}
\end{cases}
\]

这已经不是普通 causal mask。你当然可以构造完整 dense mask，但长序列下这会重新带来 $n^2$ mask 和 IO 压力。

\begin{intubox}
\textbf{为什么适合 FlexAttention：} prefix-LM 的规则很清楚，而且通常有块结构。它适合用 \texttt{mask\_mod} 表达，再转成 \texttt{BlockMask}，而不是为了这个规则专门写一个 attention kernel。
\end{intubox}

\subsection{六、最典型的需要场景 3：document mask / sample packing}

训练 LLM 时，为了提高吞吐，常常会把多个短样本 pack 到同一个长序列里。例如一个 batch 里的某条序列其实是：

\begin{center}
\begin{tabular}{p{13.2cm}}
\hline
文档 A tokens $\mid$ 文档 B tokens $\mid$ 文档 C tokens \\
\hline
\end{tabular}
\end{center}

这时如果直接用普通 causal mask，文档 B 的 token 可能会看到文档 A，文档 C 的 token 可能会看到 A/B。对语言模型训练来说，这通常是错误的数据泄漏。

正确规则应该是：

\begin{center}
\begin{tabular}{p{13.2cm}}
\hline
同一个文档内部按 causal 或 dense 规则 attend；不同文档之间不能 attend。 \\
\hline
\end{tabular}
\end{center}

概念上可以写成：

\begin{lstlisting}[language=Python]
def document_causal_mask(b, h, q_idx, kv_idx):
    same_doc = doc_id[q_idx] == doc_id[kv_idx]
    causal = kv_idx <= q_idx
    return same_doc & causal
\end{lstlisting}

\begin{summarybox}
\textbf{什么时候需要：} 只做 padding，不一定需要 FlexAttention；但如果你做 sample packing / sequence packing，并且要求不同样本之间严格隔离，\key{document mask} 就是 FlexAttention 很典型的适用场景。
\end{summarybox}

\subsection{七、最典型的需要场景 4：多模态和视频里的分区规则}

AIGC 场景里，token 往往不是单一文本序列。它可能包含：

\begin{itemize}[leftmargin=1.5em]
  \item 文本 token；
  \item 图像 patch token；
  \item 视频的时间帧 token 和空间 patch token；
  \item 条件 token、控制 token、参考图 token、全局 summary token。
\end{itemize}

这时 attention 规则可能变成：

\begin{itemize}[leftmargin=1.5em]
  \item 文本 token 可以看全部文本条件；
  \item 图像 token 只能看局部空间窗口；
  \item 视频 token 只能看邻近帧和邻近 patch；
  \item 少数全局 token 可以被所有 token 访问；
  \item 不同图、不同视频片段、不同文档之间不能互相 attend。
\end{itemize}

这些规则通常不是一个简单的 causal mask 能表达的。如果用 dense mask，mask 本身可能就很大；如果为每种组合写 kernel，工程成本又很高。

\begin{intubox}
\textbf{AIGC 使用者视角：} FlexAttention 在多模态 / 视频里最有价值的地方，不是“更快地算普通 attention”，而是让你能表达\key{文本、图像、视频、条件 token 之间的结构化可见性规则}，并尽量让底层利用这些结构化稀疏性。
\end{intubox}

\subsection{八、最典型的需要场景 5：多个 attention 规则叠加}

现实项目里麻烦的不是单个规则，而是规则组合。例如：

\begin{center}
{\small
\begin{tabular}{p{4.0cm}p{8.7cm}}
\hline
\textbf{组合} & \textbf{含义} \\
\hline
causal + sliding window & 只能看过去，而且只能看最近窗口。 \\
document mask + causal & pack 多个样本时，同文档内 causal，不同文档隔离。 \\
prefix-LM + document mask & 每个样本内部有 prefix 和 generation 区域，不同样本之间隔离。 \\
sliding window + global token & 大多数 token 看局部，少数全局 token 负责长程信息。 \\
video local window + text condition & 视频 token 看局部时空邻域，同时看文本条件。 \\
\hline
\end{tabular}
}
\end{center}

PyTorch 官方博客把这类问题称作组合爆炸：每种 mask、bias、attention setting 的组合都写一个手工 kernel，几乎不可维护。

\begin{summarybox}
\textbf{FlexAttention 的核心工程价值：} 不是为了某一个固定 mask，而是为了避免“每出现一种新 attention 变体组合，就要等别人写一个专用 kernel”。你用 PyTorch 函数描述规则，让编译器尽量生成 fused kernel。
\end{summarybox}

\subsection{九、什么时候不需要 FlexAttention？}

\begin{warnbox}
\textbf{普通场景不要过度工程化。} FlexAttention 很有用，但不是所有 attention 都应该改成 FlexAttention。
\end{warnbox}

下面这些情况通常不需要优先考虑 FlexAttention：

\begin{itemize}[leftmargin=1.5em]
  \item \textbf{普通 causal LLM attention：} 优先用 PyTorch SDPA / FlashAttention / \texttt{flash-attn} 包。
  \item \textbf{普通 dense DiT self-attention：} 如果没有复杂 mask，FlashAttention 类 backend 往往更直接。
  \item \textbf{只是简单 padding：} 先看当前框架的 SDPA / flash backend 是否已经支持，不要急着上 FlexAttention。
  \item \textbf{序列很短：} attention 不是瓶颈时，FlexAttention 的编译和 mask 构造复杂度可能不值得。
  \item \textbf{mask 每一步都高度动态且难以缓存：} \texttt{create\_block\_mask} 本身有成本，无法缓存时可能抵消收益。
  \item \textbf{mask 是随机、细粒度、没有 block 结构的稀疏：} 即使语义上 sparse，也未必能转化成好的 block-sparse 性能。
  \item \textbf{环境不适合 \texttt{torch.compile}：} 版本、硬件、动态 shape、编译时间和部署限制都可能让 FlexAttention 不合适。
  \item \textbf{decode 阶段主要瓶颈是 KV cache 管理：} 这时常常要优先看 KV cache、paged attention、batching、memory bandwidth，而不是直接套 FlexAttention。
\end{itemize}

\subsection{十、决策树：看到一个 attention 需求怎么判断？}

\begin{center}
{\small
\begin{tabular}{p{1.0cm}p{6.4cm}p{5.4cm}}
\hline
\textbf{步骤} & \textbf{问题} & \textbf{判断} \\
\hline
1 & 它是不是普通 dense 或普通 causal attention？ & 是：优先 SDPA / FlashAttention；否：继续。 \\
2 & 复杂性来自“谁能看谁”吗？ & 是：考虑 \texttt{mask\_mod}。 \\
3 & 复杂性来自“softmax 前 score 怎么改”吗？ & 是：考虑 \texttt{score\_mod}。 \\
4 & mask 是否有块结构或可跳过大块？ & 是：用 \texttt{create\_block\_mask} / \texttt{BlockMask} 更可能有性能收益。 \\
5 & mask 是否能按 shape、device、规则缓存？ & 能缓存：更适合；不能缓存：评估构造成本。 \\
6 & 现有 backend 是否已经直接支持？ & 已支持：没必要为了新而新；不支持：FlexAttention 价值上升。 \\
7 & 你是否愿意承担编译、调试和版本兼容成本？ & 愿意且收益明确：再上 FlexAttention。 \\
\hline
\end{tabular}
}
\end{center}

\begin{summarybox}
\textbf{最短判断模板：} \key{标准 attention 用 FlashAttention；复杂 mask / score 规则用 FlexAttention；复杂但没有结构或无法缓存时，要先 benchmark，不要默认一定更快。}
\end{summarybox}

\subsection{十一、FlexAttention 写法上的几个坑}

\subsubsection*{1. 用 \texttt{score\_mod} 做 mask 正确，但不一定快}

如果你写：

\begin{lstlisting}[language=Python]
def causal_score_mod(score, b, h, q_idx, kv_idx):
    return torch.where(q_idx >= kv_idx, score, -float("inf"))
\end{lstlisting}

数学语义是对的，因为被 mask 的位置 softmax 后权重为 0。但这不一定能让底层跳过对应计算。想利用 mask 稀疏性，通常应该写 \texttt{mask\_mod}，再生成 \texttt{BlockMask}。

\subsubsection*{2. \texttt{create\_block\_mask} 要缓存}

\texttt{create\_block\_mask} 不是免费的。如果你的 mask pattern 只依赖 shape、device、窗口大小、文档边界等相对稳定的信息，应该缓存它，而不是每次 forward 都重新创建。

\begin{warnbox}
\textbf{工程提示：} benchmark 时要把首次 \texttt{torch.compile} 编译时间、\texttt{BlockMask} 构造时间、warmup 后 steady-state 时间分开看。否则你可能误判 FlexAttention 的实际收益。
\end{warnbox}

\subsubsection*{3. \texttt{B=None} / \texttt{H=None} 代表广播，不是忘了写}

如果 mask pattern 不依赖 batch 或 head，可以在构造 block mask 时把 \texttt{B} 或 \texttt{H} 设成 \texttt{None}，表示这个规则在 batch/head 维度共享。这样通常更省、更容易复用。

\subsubsection*{4. 自定义函数要对编译器友好}

FlexAttention 的价值依赖 \texttt{torch.compile} 把你的规则降到 fused kernel。实践中应尽量让 \texttt{score\_mod} / \texttt{mask\_mod}：

\begin{itemize}[leftmargin=1.5em]
  \item 逻辑简单，避免复杂 Python side effect；
  \item 使用可追踪的 tensor / index 运算；
  \item 尽量避免每步变化过大的动态 shape；
  \item 先用小 shape 做正确性测试，再做性能测试。
\end{itemize}

\subsection{十二、和 Stage 1 的 FlashAttention 口径接起来}

Stage 1 里我们说过：FlashAttention 的核心是分块处理 $Q/K/V$，让局部 score tile 在片上高速内存里完成 scale、mask、online softmax 和对 $V$ 的累加，避免完整物化 $S/P$。

FlexAttention 和它的关系可以这样理解：

\begin{center}
{\small
\begin{tabular}{p{3.2cm}p{4.7cm}p{4.8cm}}
\hline
\textbf{对象} & \textbf{FlashAttention} & \textbf{FlexAttention} \\
\hline
主要目标 & 高效计算标准 dense / causal attention。 & 高效表达和计算自定义 attention 变体。 \\
核心直觉 & tiling、online softmax、fusion、少物化 $S/P$。 & 用 \texttt{score\_mod} / \texttt{mask\_mod} 描述规则，再生成 fused / block-sparse kernel。 \\
是否改变语义 & 不改变标准 attention 语义。 & 可能改变 token 可见性或 score 规则；它 exact 地执行你定义的新规则。 \\
最适合场景 & 普通 LLM causal、普通 dense self-attention。 & sliding window、prefix-LM、document mask、block-sparse、多模态复杂 mask。 \\
\hline
\end{tabular}
}
\end{center}

\begin{conceptbox}
\textbf{关键区分：} FlashAttention 是“标准 attention 的高性能实现”；FlexAttention 是“复杂 attention 规则的可编程高性能入口”。当你用 FlexAttention 改了 mask 或 score，它不再是在算普通 dense attention，而是在准确地算你定义的新 attention 语义。
\end{conceptbox}

\subsection{十三、AIGC 项目里的选型口径}

\begin{center}
{\small
\begin{tabular}{p{3.1cm}p{4.6cm}p{5.0cm}}
\hline
\textbf{场景} & \textbf{优先方案} & \textbf{什么时候转向 FlexAttention} \\
\hline
LLM 普通训练 & SDPA / FlashAttention。 & 做 sequence packing，需要 document mask；或需要 sliding window / global token 组合。 \\
LLM 长上下文 & FlashAttention + KV/cache/并行策略。 & attention pattern 从 full causal 改成局部、稀疏、prefix、文档隔离等复杂规则。 \\
DiT 图像生成 & dense attention 先用 FlashAttention 类 backend。 & patch 间只看局部窗口，或多图、多区域、多条件之间有结构化可见性规则。 \\
视频生成 & 先确认 token 数和瓶颈。 & 时间窗口、空间窗口、跨帧局部注意力、全局 token 等组合规则明显时。 \\
多模态模型 & 先看现有框架是否支持 mask。 & 文本、图像、视频、参考条件之间的 attention 规则需要自定义时。 \\
推理 decode & 先看 KV cache / paged attention / batching。 & 如果 decode/prefill 中确实有自定义 block mask 且现有 kernel 不支持，再评估。 \\
\hline
\end{tabular}
}
\end{center}

\subsection{十四、这一节的最终账本}

\begin{center}
{\small
\begin{tabular}{p{3.0cm}p{4.8cm}p{5.0cm}}
\hline
\textbf{关键词} & \textbf{核心含义} & \textbf{你要记住的工程结论} \\
\hline
\texttt{flex\_attention} & PyTorch 提供的可编程 attention API。 & 用来表达复杂 attention 变体，而不是普通 attention 的默认替代品。 \\
\texttt{score\_mod} & softmax 前修改每个 score。 & 适合 bias、soft cap、按 head/token 区域改分数。 \\
\texttt{mask\_mod} & 定义哪些 query-key pair 可见。 & 适合 sliding window、prefix-LM、document mask 等可见性规则。 \\
\texttt{BlockMask} & block 级稀疏 mask 表示。 & 让 mask 的结构化稀疏更可能变成跳过计算的性能收益。 \\
缓存 & 复用编译结果和 block mask。 & 不缓存可能让构造成本抵消性能收益。 \\
适用边界 & 规则复杂且有结构。 & 普通 dense/causal 仍优先 SDPA / FlashAttention。 \\
\hline
\end{tabular}
}
\end{center}

\begin{summarybox}
\textbf{Stage 4 最终结论：} 你需要 FlexAttention 的典型信号不是“我想让 attention 更快”，而是\key{我的 attention 规则已经不是普通 dense/causal，现有 flash backend 不直接支持；我又不想退回慢的 dense mask/eager 实现，更不想为每个规则组合手写 kernel。} 这时用 \texttt{score\_mod} 表达 score 修改，用 \texttt{mask\_mod} + \texttt{BlockMask} 表达结构化 mask，并通过 \texttt{torch.compile} 尽量获得 fused / sparse 的执行路径。
\end{summarybox}

\subsection{十五、你学完这一节后应该能立刻回答的问题}

\begin{itemize}[leftmargin=1.5em]
  \item 普通 causal LLM attention 为什么通常不需要先上 FlexAttention？
  \item FlexAttention 和 FlashAttention 的边界是什么？
  \item \texttt{score\_mod} 和 \texttt{mask\_mod} 分别解决什么问题？
  \item 为什么用 \texttt{score\_mod} 返回 $-\infty$ 可以实现 mask，但未必能获得稀疏性能收益？
  \item \texttt{BlockMask} 为什么重要？它和 dense boolean mask 的区别是什么？
  \item sliding window、prefix-LM、document mask 为什么是 FlexAttention 的典型场景？
  \item 为什么 \texttt{create\_block\_mask} 要缓存？benchmark 时为什么要区分编译时间和 steady-state 时间？
  \item 在多模态或视频生成里，什么样的 attention 规则会让你想到 FlexAttention？
\end{itemize}

% ============================================================
\subsection{参考资料}
% ============================================================

\begingroup
\small
\sloppy
\begin{itemize}[leftmargin=1.5em]
  \item PyTorch 官方博客，\emph{FlexAttention: The Flexibility of PyTorch with the Performance of FlashAttention}。用于理解 FlexAttention 的设计动机：用 PyTorch 写 attention 变体，通过 \texttt{torch.compile} 降到 fused FlashAttention-like kernel，并支持利用 mask 稀疏性。\\
  \url{https://pytorch.org/blog/flexattention/}
  \item PyTorch \texttt{torch.nn.attention.flex\_attention} 官方文档。用于确认 \texttt{flex\_attention}、\texttt{score\_mod}、\texttt{mask\_mod}、\texttt{BlockMask}、\texttt{create\_block\_mask} 等 API 语义。\\
  \url{https://docs.pytorch.org/docs/stable/nn.attention.flex_attention.html}
  \item PyTorch 源码 \texttt{flex\_attention.py}。用于校准 \texttt{flex\_attention}、\texttt{create\_block\_mask} 和 \texttt{BlockMask} 的函数签名与实现口径。\\
  \url{https://github.com/pytorch/pytorch/blob/main/torch/nn/attention/flex_attention.py}
  \item PyTorch Labs \texttt{attention-gym}：FlexAttention 示例集合，包含 masks、score mods、paged attention 等示例。\\
  \url{https://github.com/pytorch-labs/attention-gym}
  \item Dao 等，\emph{FlashAttention: Fast and Memory-Efficient Exact Attention with IO-Awareness}, arXiv:2205.14135。用于对比 FlashAttention 的标准 exact attention 高性能实现与 FlexAttention 的可编程复杂 mask 场景。\\
  \url{https://arxiv.org/abs/2205.14135}
  \item Beltagy 等，\emph{Longformer: The Long-Document Transformer}, arXiv:2004.05150。用于理解 sliding window / global attention 等长文档结构化 attention pattern 的背景。\\
  \url{https://arxiv.org/abs/2004.05150}
  \item Zaheer 等，\emph{Big Bird: Transformers for Longer Sequences}, arXiv:2007.14062。用于理解 block-sparse / global / random attention pattern 在长序列 Transformer 中的背景。\\
  \url{https://arxiv.org/abs/2007.14062}
\end{itemize}
\endgroup

